% =====================================================================
% Figure contract
% 核心结论：Native、模块化与 Multi-agent 三种范式共同收敛到边界持续迁移的 Hybrid 架构。
% 图原型：mode C / 三路范式向下汇聚到单一 hero，并由底部 summary bar 收束。
% 证据层级：hero = Hybrid 架构；支撑 = 顶部三张等权范式卡；收束 = 全宽趋势注释。
% 能不能更少？：不能；三张顶卡、三条汇聚支路、Hybrid 与趋势条均由内容规范逐项指定。
%
% Step ① 设计文档（Form A / 中等复杂度纯结构图）
% 0 参考：复用 cua-survey-fig1-overview.tex 的 XeLaTeX preamble、低饱和 palette 与描边节奏。
% 1 领域：Computer-Use Agents 中文综述；模型、系统与 agent 术语保留英文。
% 2 格式：TikZ standalone，XeLaTeX + fontspec + xeCJK；不含外部位图。
% 3 画布：W_ink 约 18.95 cm，W_print = 18.95 cm，目标印刷最小字号 6.0 pt；总宽含边框 < 20 cm。
% 4 模块：顶部三张 6.15×2.55 cm 等宽卡；中央 12.30×1.85 cm Hybrid hero；底部 18.95 cm summary。
% 5 连线：三张顶卡各自垂直下行，汇入同色 horizontal spine；唯一 arrow tip 垂直进入 Hybrid。
% 6 空间：卡间 0.25 cm；卡与 Hybrid 间 2.05 cm 作为汇聚 corridor；标签位于竖线侧方。
% 7 强调：Hybrid 使用 2.0 pt 深色描边与独有底色；顶部卡仅 0.65 pt，形成明确层级。
% 8 密度：中等档；输入无数值，不嵌入 chart/heatmap，汇聚拓扑本身承担主视觉。
% 9 ASCII 草图：
%        [ Native ]     [ 模块化 ]     [ Multi-agent ]
%             |             |               |
%             +-------------+---------------+
%                           v
%                    [ HYBRID HERO ]
%        [========== full-width convergence note ==========]
% 10 可视化决策：无可量化数据，保持普通卡片；用描边重量、留白和汇聚 spine 表达趋势。
%
% Pre-flight：P1 设计文档已写；P2 三卡按 positioning 等距构造；P3 三支路与最终 tip 已列；
% P4 N/A；P5 单一 spine corridor；P6 固定尺寸节点；P7 已读 lessons 与全局规则。
% =====================================================================

\documentclass[border=8pt]{standalone}
\usepackage{fontspec}
\usepackage{xeCJK}
\usepackage{xcolor}
\usepackage{tikz}
\usepackage{adjustbox}

\setmainfont{Helvetica Neue}
\setsansfont{Helvetica Neue}
\setCJKmainfont{Heiti SC}
\setCJKsansfont{Heiti SC}

\usetikzlibrary{arrows.meta,bending,positioning,calc}

% Low-saturation academic palette, matched to cua-survey-fig1-overview.tex.
\definecolor{ink}{HTML}{27364A}
\definecolor{muted}{HTML}{748092}
\definecolor{mainLine}{HTML}{4F667C}
\definecolor{mainFill}{HTML}{EAF0F4}
\definecolor{heroFill}{HTML}{F1F3F6}
\definecolor{heroBorder}{HTML}{3F566C}
\definecolor{blueLine}{HTML}{6D879E}
\definecolor{blueFill}{HTML}{E7EEF4}
\definecolor{purpleLine}{HTML}{817594}
\definecolor{purpleFill}{HTML}{EEEAF2}
\definecolor{goldLine}{HTML}{9B875B}
\definecolor{goldFill}{HTML}{F2EEE3}
\definecolor{tealLine}{HTML}{648D8A}
\definecolor{tealFill}{HTML}{E5F0EE}

\tikzset{
  every node/.style={
    font=\sffamily\fontsize{8.2}{10.0}\selectfont,
    text=ink
  },
  figuretitle/.style={
    font=\sffamily\fontsize{12.0}{14.2}\selectfont\bfseries,
    text=heroBorder,
    align=center
  },
  topcard/.style={
    rectangle,
    rounded corners=5pt,
    line width=0.65pt,
    minimum width=6.15cm,
    minimum height=2.55cm,
    text width=5.72cm,
    align=center,
    inner xsep=6pt,
    inner ysep=7pt
  },
  hybridcard/.style={
    rectangle,
    rounded corners=6pt,
    draw=heroBorder,
    fill=purpleFill!68!white,
    line width=2.0pt,
    minimum width=12.30cm,
    minimum height=1.85cm,
    text width=11.60cm,
    align=center,
    inner xsep=9pt,
    inner ysep=8pt
  },
  summarybar/.style={
    rectangle,
    rounded corners=4pt,
    draw=muted!55,
    fill=mainFill!58,
    line width=0.55pt,
    minimum width=18.95cm,
    minimum height=1.22cm,
    text width=18.20cm,
    align=center,
    inner xsep=7pt,
    inner ysep=6pt,
    font=\sffamily\fontsize{7.35}{9.0}\selectfont
  },
  branch/.style={
    line width=0.9pt,
    color=mainLine
  },
  arrow short/.style={
    -{Stealth[length=3pt 1.0,width'=0pt 0.5,sep=0pt 0.3]},
    line width=0.9pt,
    shorten >=1pt,
    shorten <=0pt,
    color=mainLine
  },
  connectorlabel/.style={
    font=\sffamily\fontsize{7.1}{8.6}\selectfont,
    text=mainLine,
    inner sep=0pt
  }
}

\newcommand{\mechanism}[1]{%
  \adjustbox{max width=5.55cm}{\fontsize{7.0}{8.4}\selectfont\strut #1}%
}

\begin{document}
\sffamily
\begin{tikzpicture}

\node[figuretitle] (title) {模型与系统范式(§6.1–6.6)};

% Three visually symmetric top-row cards; each contains exactly three text lines.
\node[
  topcard,
  below=0.48cm of title,
  fill=goldFill,
  draw=goldLine
] (modular) {
  {\fontsize{9.2}{11.0}\selectfont\bfseries 模块化系统 §6.4}\\[5pt]
  \mechanism{grounder·planner·memory·critic 专用模块分工}\\[5pt]
  {\color{muted}\fontsize{7.6}{9.0}\selectfont 例:Agent S2}
};

\node[
  topcard,
  left=0.25cm of modular,
  fill=blueFill,
  draw=blueLine
] (native) {
  {\fontsize{9.2}{11.0}\selectfont\bfseries Native 端到端 §6.3}\\[5pt]
  \mechanism{perception·grounding·reasoning·action 收进同一可训练模型}\\[5pt]
  {\color{muted}\fontsize{7.6}{9.0}\selectfont 例:UI-TARS·ScaleCUA}
};

\node[
  topcard,
  right=0.25cm of modular,
  fill=tealFill,
  draw=tealLine
] (multi) {
  {\fontsize{9.2}{11.0}\selectfont\bfseries Multi-agent §6.5}\\[5pt]
  \mechanism{Orchestrator 编排多个执行 agent}\\[5pt]
  {\color{muted}\fontsize{7.6}{9.0}\selectfont 例:CoAct-1}
};

\node[hybridcard,below=2.05cm of modular] (hybrid) {
  {\fontsize{10.0}{12.0}\selectfont\bfseries Hybrid 架构 §6.6 — 模型与系统的边界迁移}\\[5pt]
  {\color{muted}\fontsize{8.0}{9.6}\selectfont 例:UI-TARS-2·StateAct}
};

% Canonical fan-in: three tip-free branches -> one continuous spine -> one target tip.
\coordinate (spineCenter) at ($(modular.south)!0.66!(hybrid.north)$);
\coordinate (nativeJoin) at (native.south |- spineCenter);
\coordinate (modularJoin) at (modular.south |- spineCenter);
\coordinate (multiJoin) at (multi.south |- spineCenter);

\draw[branch] (native.south) -- (nativeJoin);
\draw[branch] (modular.south) -- (modularJoin);
\draw[branch] (multi.south) -- (multiJoin);
\draw[branch] (nativeJoin) -- (multiJoin);
\draw[arrow short] (modularJoin) -- (hybrid.north);

% Horizontal labels sit beside, never on top of, their vertical branches.
\node[connectorlabel,anchor=east,xshift=-0.20cm]
  at ($(native.south)!0.45!(nativeJoin)$) {GUI 与 SDK};
\node[connectorlabel,anchor=west,xshift=0.20cm]
  at ($(modular.south)!0.40!(modularJoin)$) {planner 与 specialist grounder};

\node[summarybar,below=0.52cm of hybrid] (summary) {
  \mbox{收敛趋势:native 化不消除系统设计,只把系统边界从显式模块接口迁移到训练数据、context policy 与 action schema(§6.3、§11.3)}
};

\end{tikzpicture}
\end{document}
